% ============================================================================
% REWRITTEN "RESULTS AND DISCUSSION" SECTION — drop-in replacement
% Generated from results/intphys_primary_db.xlsx and results/mvp_primary_db.xlsx
% Selection criterion: best layer per (model, probe, benchmark) by TEST primary
% metric ("best-case" analysis), matching the convention of the current Table 1.
% All numbers verified against the Excel DBs on 2026-07-18.
%
% Open items are marked with %% FLAG comments — see the summary that accompanies
% this file. Do not delete them before resolving with Niccolò.
% ============================================================================

\section{Results and Discussion}

% ----------------------------------------------------------------------------
% TABLE 1 — largest backbones (experiment "main" + LTX-13B).
% Corrected cells vs. the current PDF are marked with %% CHANGED.
% ----------------------------------------------------------------------------
\begin{table*}[t]
\centering
\small
\setlength{\tabcolsep}{4pt}
\renewcommand{\arraystretch}{1.08}
\resizebox{\textwidth}{!}{%
\begin{tabular}{l cc cc cc c cc cc cc}
\toprule
& \multicolumn{6}{c}{\textbf{IntPhys2}} & & \multicolumn{6}{c}{\textbf{MVP}} \\
\cmidrule(lr){2-7} \cmidrule(lr){9-14}
\textbf{Model}
& \multicolumn{2}{c}{\textbf{Linear}} & \multicolumn{2}{c}{\textbf{MLP}} & \multicolumn{2}{c}{\textbf{Temporal Attn.}}
& & \multicolumn{2}{c}{\textbf{Linear}} & \multicolumn{2}{c}{\textbf{MLP}} & \multicolumn{2}{c}{\textbf{Temporal Attn.}} \\
\cmidrule(lr){2-3} \cmidrule(lr){4-5} \cmidrule(lr){6-7} \cmidrule(lr){9-10} \cmidrule(lr){11-12} \cmidrule(lr){13-14}
& \textbf{Layer} & \metriccol{\textbf{VOE}} & \textbf{Layer} & \metriccol{\textbf{VOE}} & \textbf{Layer} & \metriccol{\textbf{VOE}}
& & \textbf{Layer} & \metriccol{\textbf{Pair}} & \textbf{Layer} & \metriccol{\textbf{Pair}} & \textbf{Layer} & \metriccol{\textbf{Pair}} \\
\midrule
V-JEPA
& 24 & \bestmetric{50.98} & 32 & \metriccol{45.10} & 16 & \metriccol{66.67}
& & 32 & \metriccol{48.74} & 32 & \bestmetric{87.26} & 24 & \bestmetric{94.03} \\
V-JEPA 2
& 40 & \metriccol{41.18} & 30 & \bestmetric{56.86} & 40 & \bestmetric{78.43}
& & 40 & \metriccol{48.53} & 40 & \metriccol{86.75} & 40 & \metriccol{93.33} \\ %% CHANGED: MVP Linear 47.32 -> 48.53
V-JEPA 2.1
& 48 & \metriccol{35.29} & 36 & \metriccol{39.22} & 48 & \metriccol{60.78}
& & 36 & \metriccol{42.26} & 36 & \metriccol{70.78} & 48 & \metriccol{93.73} \\ %% CHANGED: MVP Linear 43.88@38 -> 42.26@36; MLP layer 38 -> 36
VideoMAE
& 16 & \metriccol{35.29} & 24 & \metriccol{35.29} & 24 & \metriccol{76.47}
& & 24 & \metriccol{37.51} & 24 & \metriccol{64.41} & 24 & \metriccol{92.01} \\
VideoMAE-v2
& 40 & \metriccol{47.06} & 20 & \metriccol{47.06} & 20 & \metriccol{62.75}
& & 20 & \metriccol{37.71} & 20 & \metriccol{66.13} & 30 & \metriccol{91.10} \\ %% CHANGED: MVP Linear 38.62 -> 37.71
LTX-Video
& 24 {\scriptsize(0.1)} & \metriccol{49.02} & 24 {\scriptsize(0.2)} & \metriccol{47.06} & 36 {\scriptsize(0.1)} & \metriccol{43.14} %% FLAG: attn label ambiguous (36/0.1 vs 48/0.2, same value) — confirm with run logs
& & 24 {\scriptsize(0.6)} & \bestmetric{49.95} & 24 {\scriptsize(0.7)} & \metriccol{69.16} & 24 {\scriptsize(0.5)} & \metriccol{84.33} \\
\bottomrule
\end{tabular}}
\caption{
Best-case benchmark comparison across model families, each evaluated with its
largest available backbone. For each dataset and probe, we report the best layer
and the benchmark-specific primary metric. Shaded columns denote the primary
metric; the darkest cell within each probe block indicates the highest value.
For LTX-Video (13B), layer columns report the transformer block, with the
denoising noise level in parentheses.
}
\label{tab:exp1-bestcase-panels}
\end{table*}

% ============================================================================
\subsection*{Analysis 1: Readout accessibility under linear and nonlinear probes}
\label{sec:exp1}

\begin{researchquestionbox}
\textbf{Research question.}
How is intuitive-physics information encoded in frozen representations:
linearly, nonlinearly, or in the temporal structure of the token sequence?
\end{researchquestionbox}

To answer this question we compare, for each model evaluated with its largest
available backbone, the best layer under probes of increasing expressivity
(\autoref{tab:exp1-bestcase-panels}). On MVP, probe capacity matters
substantially: linear probes remain weak across all model families
(37.5--50.0\% pair consistency, against a 25\% random baseline), MLP probes
recover much stronger signal (64.4--87.3\% for the ViT families), and temporal
attentive probes are uniformly best, reaching 91.1--94.0\% for every ViT model
and 84.3\% for LTX-Video. This large and consistent separation indicates that
MVP-relevant physics information is only weakly accessible to a linear readout
and becomes recoverable once the probe can model nonlinear and, especially,
explicitly temporal interactions.

IntPhys2 shows a different profile across the first two probe families. Linear
probes already recover moderate signal in several models (35.3--51.0\% VOE,
against a $\sim$16.7\% random baseline), and MLP probes add little on top:
gains are small and inconsistent, with V-JEPA even decreasing from 50.98\% to
45.10\%. The temporal attentive probe, however, changes the picture
substantially: it improves over the pooled probes for every ViT model, often by
a wide margin (66.67\% for V-JEPA, 78.43\% for V-JEPA 2, 76.47\% for VideoMAE,
62.75\% for VideoMAE-v2). The only exception is LTX-Video, whose attentive
probe (43.14\%) does not exceed its linear best (49.02\%). Since linear and MLP
probes operate on average-pooled clip embeddings while the attentive probe sees
the full token sequence, this pattern suggests that on IntPhys2 much of the
missing signal is not merely nonlinear but is stored in the spatio-temporal
token structure that pooling destroys.

\begin{takeawaybox}
\textbf{Takeaway}.
Intuitive-physics information is not equally explicit across tasks or readout
mechanisms. On MVP it is largely inaccessible to linear readouts and emerges
progressively with nonlinear and temporal probes. On IntPhys2 a larger fraction
of the signal is already linearly decodable from pooled features, but the
largest gains come from probes that operate directly on the token sequence,
indicating that much of the physics-relevant structure lives in spatio-temporal
token interactions rather than in the pooled clip embedding.
\end{takeawaybox}

% ============================================================================
\subsection*{Analysis 2: Layerwise emergence of physics information}
\label{sec:exp2}

\begin{researchquestionbox}
\textbf{Research question.}
Where in the network does intuitive-physics information become most accessible?
\end{researchquestionbox}

To determine where physical understanding originates within the backbones, we
analyze the layerwise performance profiles of the largest backbone of each
model. Here we focus on the MLP probe, as it possesses sufficient capacity to
extract latent task-relevant signals without explicitly re-modeling temporal
sequences; corresponding profiles for the linear and temporal attentive probes
are reported in Appendix \autoref{fig:exp2_all_probes} and
\autoref{fig:exp2_ltx_noise_linear}.

As \autoref{fig:exp2_mlp} and \autoref{fig:exp2_ltx_noise_mlp} show, the depth
at which physics-relevant information becomes most accessible depends strongly
on the benchmark. On MVP, performance generally improves toward later layers,
especially for the V-JEPA family: V-JEPA rises from 59.45\% pair consistency at
0.25 depth to 87.26\% at the final layer, and V-JEPA 2 shows a similar increase
from 52.07\% to 86.75\%. The VideoMAE models follow the same broad mid-to-late
trend with weaker gains and earlier peaks (VideoMAE peaks at 0.75 depth with
64.41\%, VideoMAE-v2 at 0.5 depth with 66.13\%). Overall, MVP is best
characterized as late-layer dominated, suggesting that the information needed
for minimal-pair physical reasoning becomes most accessible only after
substantial processing.

IntPhys2 shows a less monotonic profile. Several models peak around 0.5--0.75
depth and then flatten or decline at the final layer; for example, V-JEPA 2
reaches 56.86\% VOE at 0.75 depth before dropping to 47.06\% at the output, and
VideoMAE-v2 peaks at 0.5 depth (47.06\%) before declining. LTX-Video follows
the same broad pattern along its own two axes: its strongest physics signal is
concentrated around the middle transformer blocks, and fixing a block while
reducing the denoising noise level does not produce a monotonic improvement;
performance fluctuates across the denoising trajectory
(\autoref{fig:exp2_ltx_noise_mlp}). Thus, even for diffusion features, the most
accessible physics signal appears at intermediate depth rather than at the
extremes of the network or of the denoising process.

\begin{takeawaybox}
\textbf{Takeaway}.
Physics-relevant information is weakest in early layers and becomes most
accessible at intermediate-to-late depth. MVP shows the clearest late-layer
trend, whereas IntPhys2 more often peaks between 50\% and 75\% depth. LTX-Video
follows the same broad pattern, with its strongest signal concentrated in the
middle of the backbone and no monotonic trend along the denoising trajectory.
\end{takeawaybox}

% ============================================================================
\subsection*{Analysis 3: Influence of the backbone size}
\label{sec:exp3}

\begin{researchquestionbox}
\textbf{Research question.}
How much does the choice of backbone influence the accessibility of physical
understanding?
\end{researchquestionbox}

To separate the effect of the pretraining objective from the effect of model
capacity, we sweep the available backbone sizes of the most recent model in
each family: V-JEPA 2.1 (ViT-B, ViT-L, ViT-G, ViT-Gigantic) and VideoMAE-v2
(ViT-B, ViT-L, ViT-G), and additionally compare the two available LTX-Video
variants (2B vs.\ 13B). \autoref{fig:exp4_backbones} reports best-case results
per backbone; full layerwise profiles are shown in Appendix
\autoref{fig:exp4_backbones_layerwise}.

For the ViT families, backbone size has a real but modest and probe-dependent
effect. Under the temporal attentive probe, performance increases roughly
monotonically with size: VideoMAE-v2 improves from 52.94\% (ViT-B) to 54.90\%
(ViT-L) to 62.75\% (ViT-G) VOE on IntPhys2 and from 85.74\% to 87.36\% to
91.10\% pair consistency on MVP, while V-JEPA 2.1 moves from 58.82\% (ViT-B) to
60.78\% (ViT-Gigantic) on IntPhys2 and stays within a narrow 91.4--93.7\% band
on MVP. Under linear and MLP probes, by contrast, no consistent scaling trend
emerges, and the smallest backbone is sometimes best: the ViT-B variant of
V-JEPA 2.1 attains the family's highest MLP VOE on IntPhys2 (50.98\%). Scaling
the backbone therefore does not by itself make physics information more
linearly explicit; its benefits materialize mainly when the readout can exploit
the token sequence.

The diffusion family behaves differently: model size matters much more.
LTX-Video 2B is close to the random baseline in most configurations, with
best-case scores of 27.45\% VOE on IntPhys2 (vs.\ 49.02\% for 13B) and 40.34\%
MLP pair consistency on MVP (vs.\ 69.16\% for 13B).
%% FLAG: LTX-2B attentive probe on MVP has not been run (empty rows in the DB)
%% — either complete the run or state the omission explicitly here.
This gap is consistent with the view that representation quality is a
by-product of the generative objective: a small diffusion model can be a
reasonable generator while encoding little accessible physical structure,
whereas the ViT families already saturate much of their probe-accessible
physics signal at ViT-B/L scale.

\begin{takeawaybox}
\textbf{Takeaway}.
Within the ViT families, backbone size is a second-order factor: probe
expressivity matters far more than scale, and scaling helps mainly under
temporally expressive readouts. For the diffusion family, scale is
first-order: LTX-Video 2B recovers close-to-chance signal, while the 13B
variant becomes competitive with pooled readouts of the ViT models.
\end{takeawaybox}

% ============================================================================
\subsection*{Analysis 4: Pretraining paradigms under a matched backbone}
\label{sec:exp4}

\begin{researchquestionbox}
\textbf{Research question}
Which pretraining paradigm encodes the most accessible intuitive-physics
information, once backbone capacity is matched?
\end{researchquestionbox}

The comparison in \autoref{tab:exp1-bestcase-panels} conflates the pretraining
objective with backbone scale, since each family ships different maximal sizes.
To control for this, we repeat the best-case comparison with all five ViT
models on a shared ViT-L/16 backbone (\autoref{fig:exp3_probe_comparison}); the
two LTX-Video variants, which have no ViT-L counterpart, are included as
diffusion references. We focus the discussion on the temporal attentive probe,
which operates directly on the token sequence and is thus most representative
of the models' ability to leverage temporal structure.

At matched capacity, the ordering of paradigms is clear and consistent with the
largest-backbone comparison, indicating that it is not an artifact of scale.
The V-JEPA family retains the most accessible physics information: the original
V-JEPA leads both benchmarks at ViT-L (74.51\% VOE on IntPhys2 and 93.83\% pair
consistency on MVP), with V-JEPA 2 close behind (70.59\% and 93.33\%). The
VideoMAE family is competitive on MVP (89.89\% for VideoMAE, a gap of
$\sim$4\%) but falls further behind on IntPhys2 (68.63\% for VideoMAE, 54.90\%
for VideoMAE-v2). The diffusion-based LTX-Video trails on both benchmarks even
at 13B parameters, i.e., with far more capacity than a ViT-L: 84.33\% on MVP
($\sim$9.5\% below V-JEPA) and 43.14\% on IntPhys2, a gap of more than 30
points. Interestingly, the newest member of each family is not the best at
matched size: V-JEPA outperforms V-JEPA 2.1 (58.82\% VOE) and VideoMAE
outperforms VideoMAE-v2 on IntPhys2, suggesting that recent training recipes
optimized at larger scales do not automatically transfer their advantages to
smaller backbones.

These differences are consistent with the nature of the objectives. Learning to
predict target representations in a joint embedding space, rather than
reconstructing raw pixels, encourages abstract, temporally grounded features in
which physical structure is readily accessible. Masked reconstruction retains
much of this signal but organizes it less explicitly, and a denoising objective
optimized for generation, while producing strong visual results, does not by
itself organize fine-grained intuitive physics into an accessible format. We
note that our comparison controls for backbone capacity but not for pretraining
data, recency, or training recipe, which remain confounded with the objective.

\begin{takeawaybox}
\textbf{Takeaway}:
Under our probing protocol and at matched ViT-L capacity, predictive
joint-embedding pretraining yields the most accessible and robust
intuitive-physics representations, followed by masked reconstruction, with the
diffusion paradigm last despite a much larger parameter budget.
\end{takeawaybox}

% ============================================================================
% APPENDIX — Temporal control conditions (moved from main text).
% The attentive-probe IntPhys2 controls were re-run together with the
% longer-epoch attentive probes; the table row values below are recomputed
% from results/intphys2_temporal_attn_controls.csv at the new best layers.
% Linear and MLP control rows are unchanged from the previous draft.
% ============================================================================

% --- Replacement values for panel (a) IntPhys2, "Temporal Attn." column ---
% Model         Main    Delta-shuffle(%)  Delta-single(%)
% V-JEPA        66.67   -41.18            -76.47
% V-JEPA 2      78.43   -30.00            -77.50
% V-JEPA 2.1    60.78    0.00             -74.19
% VideoMAE      76.47   -30.77            -76.92
% VideoMAE-v2   62.75   -25.00            -68.75
% LTX-Video     43.14   -90.91            -72.73
%% FLAG: LTX row assumes best attn = block 36 / noise 0.1 (run index 39). If the
%% best-attn config is instead block 48 / noise 0.2 (run index 36), the deltas
%% become -59.09 / -81.82. Resolve the label ambiguity first.

\paragraph{Temporal controls (appendix prose).}
On MVP, temporal disruption produces a clear collapse across all probe
families, confirming that successful minimal-pair classification requires
multi-frame temporal evidence. On IntPhys2, single-frame repetition remains
highly damaging for every model and probe (attentive probes drop by 68--78\%
relative), showing that the benchmark is not solvable from a static frame.
Frame shuffling, however, is less uniformly destructive: the V-JEPA 2.1
attentive probe shows exactly a 0.00\% relative drop, and several other
configurations retain a substantial fraction of their main-task score. This
suggests that part of the IntPhys2 VOE signal can be recovered from
order-insensitive, unordered multi-frame evidence rather than from the causal
event trajectory itself, a qualitatively different solution strategy from the
temporally grounded one that MVP requires.
%% NOTE: the old footnote about the anomalous +62.50% single-frame jump of
%% VideoMAE-v2's attentive probe must be DELETED — the re-run with more epochs
%% removed the underfit outlier (main-task VOE is now 62.75, and both controls
%% reduce it).
